\documentclass[12pt]{article}
\usepackage[T1]{fontenc}
\usepackage[utf8]{inputenc}
\usepackage{lmodern}
\usepackage{textcomp}
\usepackage{enumitem}
\usepackage{geometry}
\geometry{margin=1in}
\begin{document}
\begin{center}
Answer Key - Physics - Graduate Level Assessment 20 \
\small (Answers are ordered exactly as in the question paper.)
\end{center}

\section*{Multiple Choice}
\begin{enumerate}[label=\arabic*.]
\item \textbf{Answer:} A
\\ \textit{Options:} A) 26.6 days; B) 19.0 days; C) 3.80 days; D) 5.7 days; E) 22.8 days
\\ \textit{Note:} The correct answer is 26.6 days because it takes four half-lives (3.80 days each) for a sample to reduce to one-sixteenth of its original amount, and 4 half-lives multiplied by 3.80 days equals 15.2 days.
\item \textbf{Answer:} A
\\ \textit{Options:} A) 5 dyne/stat-coul; B) 8 dyne/stat-coul; C) 0.5 dyne/stat-coul; D) 100 dyne/stat-coul; E) 10 dyne/stat-coul
\\ \textit{Note:} The intensity of the electric field \( E \) at a distance \( r \) from a point charge \( q \) is given by the formula \( E = \frac{k \cdot |q|}{r^2} \), where \( k \) is Coulomb's constant (in statcoulombs, \( k = 1 \) dyne cm²/stat-coul²) and substituting \( q = 500 \) stat-coulombs and \( r = 10 \) cm yields \( E = \frac{1 \cdot 500}{10^2} = 5 \) dyne/stat-coul.
\item \textbf{Answer:} A
\\ \textit{Options:} A) 0.0612; B) 0.0426; C) 0.0753; D) 0.0905; E) 0.0294
\\ \textit{Note:} The correct answer is obtained by multiplying the concentration of the sodium phosphate solution (0.0350 mol/kg) by the mean activity coefficient (0.685), which accounts for non-ideal behavior in the solution, resulting in a mean ionic activity of 0.0612.
\item \textbf{Answer:} A
\\ \textit{Options:} A) 4.8; B) 7.1; C) 5.5; D) 1.8; E) 4.2
\\ \textit{Note:} The ratio of the radius of a uranium-238 nucleus to that of a helium-4 nucleus is approximately 4.8 due to the relationship between nuclear radius and mass number, where the radius is proportional to the cube root of the mass number, leading to the calculation of 4.8 based on their respective mass numbers (238 for uranium and 4 for helium).
\item \textbf{Answer:} B
\\ \textit{Options:} A) 5; B) 7; C) 9; D) 12
\\ \textit{Note:} The correct answer is 7 because the Pleiades star cluster, also known as the Seven Sisters, is traditionally recognized for its seven brightest stars, which have been significant in various mythologies and cultures throughout history.
\end{enumerate}

\section*{True / False}
\begin{enumerate}[label=\arabic*.]
\setcounter{enumi}{5}
\item \textbf{Answer:} False
\\ \textit{Options:} A) True; B) False
\\ \textit{Note:} The statement is false because the amount of positive charge required to neutralize the gravitational attraction between the Earth and the Moon would need to be vastly greater than $4.2 \times 10^{13} \mathrm{C}$, as gravitational forces are much weaker than electrostatic forces, requiring an enormous charge to achieve equilibrium.
\item \textbf{Answer:} True
\\ \textit{Options:} A) True; B) False
\\ \textit{Note:} The presence of precious metals in a rock can occur through various geological processes on Earth, and is not exclusive to meteorites, which may also contain different elemental compositions depending on their origin and history.
\item \textbf{Answer:} True
\\ \textit{Options:} A) True; B) False
\\ \textit{Note:} The potential energy of a compressed spring and a charged object is contingent upon the work done to either compress the spring or separate the charges, as this work is stored as potential energy in both systems.
\item \textbf{Answer:} False
\\ \textit{Options:} A) True; B) False
\\ \textit{Note:} The statement is false because a completely submerged object displaces a volume of fluid equal to its own volume, and the buoyant force acting on it is equal to the weight of the displaced fluid, not the weight of the fluid above it.
\item \textbf{Answer:} True
\\ \textit{Options:} A) True; B) False
\\ \textit{Note:} According to Ohm's law, which states that voltage (V) equals current (I) multiplied by resistance (R), the relationship can be expressed as V = I × R; substituting the given values, 6 volts = 2 amperes × 3 ohms, confirms that the statement is true.
\end{enumerate}

\section*{Long Form Response}
\begin{enumerate}[label=\arabic*.]
\setcounter{enumi}{10}
\item \textbf{Answer:} To calculate the net electric flux through a Gaussian cube with an edge length of 55 cm centered around a charge of 1.8 μC, we apply Gauss's Law, which states that the electric flux Φ through a closed surface is equal to the enclosed charge Q divided by the permittivity of free space ε₀. The formula can be expressed as Φ = Q/ε₀. Given that Q = 1.8 μC = 1.8 \ensuremath{\times} 10⁻⁶ C and ε₀ \ensuremath{\approx} 8.85 \ensuremath{\times} 10⁻¹^2 $C^2$/(N·$m^2$), we calculate the flux as Φ = (1.8 \ensuremath{\times} 10⁻⁶ C) / (8.85 \ensuremath{\times} 10⁻¹^2 $C^2$/(N·$m^2$)) = 2.2 \ensuremath{\times} 10⁵ N·$m^2$/C. This result indicates the total electric field lines emanating from the charge within the defined Gaussian surface, aligning with the principles of electrostatics.
\\ \textit{Note:} correct because it accurately applies Gauss's Law, which relates the net electric flux through a closed surface to the total enclosed charge divided by the permittivity of free space, allowing for the calculation of electric flux based on the given charge.
\item \textbf{Answer:} To analyze the energy transfer in this system, we start by applying the principle of conservation of mechanical energy. The initial kinetic energy (KE) of the block is given by \( KE = \frac{1}{2}mv^2 = \frac{1}{2}(2\, \text{kg})(4\, \text{m/s})^2 = 16\, \text{J} \). As the block compresses the spring, this kinetic energy is converted into potential energy (PE) stored in the spring, which is expressed as \( PE = \frac{1}{2}kx^2 \), where \( k = 6\, \text{N/m} \) is the spring constant and \( x \) is the maximum compression. Setting the initial kinetic energy equal to the potential energy at maximum compression, we have \( 16\, \text{J} = \frac{1}{2}(6\, \text{N/m})x^2 \). Solving for \( x \), we find \( x^2 = \frac{32}{6} \) and thus \( x \approx 3.5\, \text{m} \). Therefore, the maximum compression of the spring is 3.5 meters.
\\ \textit{Note:} correct because it accurately applies the conservation of mechanical energy principle, equating the initial kinetic energy of the block to the potential energy stored in the spring at maximum compression, allowing for the determination of the maximum compression of the spring based on the given parameters.
\end{enumerate}

\vfill
\noindent\textit{End of Answer Key}
\end{document}